\section{Testes dos algoritmos de controle}

Esta seção descreve os ambientes de simulação, os parâmetros numéricos e os modelos de medição adotados para a execução das simulações numéricas.
Métricas como desempenho, estabilidade ou robustez dos controladores não são discutidos neste momento, pois serão introduzidos no capítulo de simulações.

\subsection{MATLAB/Simulink}
O MATLAB/Simulink foi a principal plataforma de simulação em tempo contínuo. 
A modelagem matemática do problema orbital e de atitude foi implementado de forma modular, para integrar cada etapa da modelagem matemática apresentada neste capítulo.
A integração numérica foi realizada pelo método ode45 (Runge--Kutta) de quarta/quinta ordem com passo adaptativo. 

\subsection{Ambiente virtual 3D no Gazebo}
Além do MATLAB/Simulink, foi utilizado o simulador Gazebo, que permite a modelagem tridimensional e a simulação física em tempo real. 
Nesse ambiente, o modelo CAD da espaçonave visitante foi importada com suas respectivas propriedades geométricas e inerciais. 
O Gazebo serviu como plataforma para validar o controle em um cenário mais próximo da operação real, possuindo interação direta entre o \emph{chaser} e o ambiente.

\subsection{Condições iniciais}
O estado inicial do \emph{chaser} foi ajustado de acordo com o cenário de aproximação considerado, que abrangeu somente órbita coplanar.

\subsection{Subsistemas}
No subsistema orbital, a aproximação final é realizada ao longo do eixo radial (\emph{R-bar}) e a transferência inicial entre órbitas circulares é obtida por meio da transferência de Hohmann.  
Para o subsistema de atitude são analisadas as equações de movimento da espaçonave com base flutuante sujeita à ação dos torques de controle, bem como as equações de movimento das juntas do manipulador composto por quatro articulações revolutas e uma prismática (configuração 4R1P), responsáveis pelo alinhamento e extensão da ferramenta de captura.

\subsection{Modelagem de sensores}
A realimentação das variáveis de estado foi realizada por meio de sensores virtuais (giroscópios e acelerômetros para estimar velocidades angulares e acelerações, sensores de navegação para determinar a atitude e a posição relativa das espaçonaves) modelados nos ambientes de simulação. 

Com os ambientes, parâmetros e modelos definidos, o Capítulo 3 apresenta os resultados dos testes, as métricas adotadas e a análise comparativa de desempenho.

\subsection{Critérios de Desempenho}
\label{sec:criterios_desempenho}

Os critérios de desempenho permitem quantificar indicares de qualidade do modelo implementado. As métricas utilizadas são:
\begin{itemize}
  \item erro (desejado e atual) de posição e de ângulo;
  \item tempo de estabilização;
  \item esforço do controle, onde é avaliado a existência de torques ou de forças irreais;
  \item comportamento em regime permanente;
  \item saturação dos atuadores, verificando se os torques ou forças exigidos permanecem dentro dos limites atribuídos ao sistema;
\end{itemize}

As métricas completas podem ser visualizadas no Apêndice \eqref{ApendiceA}.




% \input{2 Formulacao do problema/2.9 Testes de algoritmos/2.9.1 Ambiente de simulacao MatlabSimulink.tex}
% \input{2 Formulacao do problema/2.9 Testes de algoritmos/2.9.2 Ambiente virtual 3D Gazebo.tex}
% \input{2 Formulacao do problema/2.9 Testes de algoritmos/2.9.3 Parametros de entrada e CI.tex}
% \input{2 Formulacao do problema/2.9 Testes de algoritmos/2.9.4 Modelagem sensores e ruidos.tex}



